%==============================================================================
%  Master Data Entity Resolution -- proposal
%
%  Compiles with pdflatex on a stock TeX Live / Overleaf install; every package
%  below ships in the base distribution, so there is nothing to install.
%
%      pdflatex mder_proposal.tex   (run twice, for the table of contents)
%
%  Edit \orgname and \authorname below before circulating.
%==============================================================================
\documentclass[11pt,a4paper]{article}

\usepackage[T1]{fontenc}
\usepackage[utf8]{inputenc}
\usepackage{lmodern}
\usepackage[margin=2.2cm]{geometry}
\usepackage{booktabs}
\usepackage{array}
\usepackage{xcolor}
\usepackage{titlesec}
\usepackage{enumitem}
\usepackage{caption}
\usepackage[hidelinks]{hyperref}

%------------------------------------------------------------------ house style
\definecolor{accent}{HTML}{003DA5}
\definecolor{rulegrey}{HTML}{C8CDD4}
\definecolor{panelbg}{HTML}{F2F5F9}

\titleformat{\section}{\large\bfseries\color{accent}}{\thesection}{0.6em}{}
\titleformat{\subsection}{\normalsize\bfseries}{\thesubsection}{0.6em}{}
\titlespacing*{\section}{0pt}{1.35em}{0.5em}
\titlespacing*{\subsection}{0pt}{0.9em}{0.35em}

\setlist[itemize]{leftmargin=1.2em,itemsep=0.25em,topsep=0.3em}
\setlength{\parskip}{0.45em}
\setlength{\parindent}{0pt}
\renewcommand{\arraystretch}{1.15}

% a light call-out panel
\newcommand{\panel}[1]{%
  \par\vspace{0.4em}%
  \noindent\colorbox{panelbg}{%
    \begin{minipage}{\dimexpr\linewidth-2\fboxsep}%
      \vspace{0.35em}#1\vspace{0.35em}%
    \end{minipage}}%
  \par\vspace{0.5em}%
}

%--------------------------------------------------------------------- metadata
\newcommand{\orgname}{[Organisation]}
\newcommand{\authorname}{Jacob Bethell}

\begin{document}

%================================================================== title block
\begin{center}
  {\LARGE\bfseries\color{accent} Master Data Entity Resolution}\\[0.35em]
  {\large Automated resolution of supplier product submissions\\
   against a master catalogue}\\[0.9em]
  {\normalsize \authorname\quad\textbullet\quad Working prototype and evaluation}\\[0.2em]
  {\small\color{black!55} Prepared for \orgname\ --- \today}
\end{center}

\vspace{0.4em}
{\color{rulegrey}\hrule height 1.2pt}
\vspace{1.0em}

%============================================================ executive summary
\section*{Summary}

Every product a supplier sends us arrives described three different ways: what
is printed on the packaging, what the supplier typed into our portal, and what
our master catalogue already holds. Deciding whether those three describe the
same product is currently manual work, and getting it wrong in either direction
is expensive --- a duplicate record fragments sales reporting and stock, while a
wrong merge silently corrupts a record that many downstream systems trust.

I have built a working eight-stage pipeline that makes that decision
automatically, attaches a calibrated confidence figure to it, and routes each
submission to one of three outcomes: merge without review, queue for a human, or
reject. It is evaluated on three independent datasets, including 2{,}256 real
grocery products and 140 real packaging photographs.

\panel{%
\textbf{What I am asking for:} a time-boxed pilot on a real supplier feed --- roughly
\textbf{six weeks}, one engineer, plus a data steward for two days a week to label
outcomes. The prototype already exists and is tested; what it has never seen is
\emph{our} catalogue and \emph{our} suppliers. The pilot answers one question:
what share of submissions can be resolved without a human, at a wrong-merge rate
we are willing to accept?}

%=================================================================== the problem
\section{The problem}

A supplier submission has to be matched to exactly one catalogue record before
anything downstream can use it. The difficulty is that the three descriptions
rarely agree:

\begin{itemize}
  \item the packaging says \textsf{Mars}, the portal submission says
        \textsf{Mars Inc}, the catalogue says \textsf{Mars};
  \item the pack says \textsf{51g}, the submission says \textsf{0.051kg};
  \item and roughly one supplier submission in three contains at least one
        outright keying error.
\end{itemize}

Two distinct failure modes follow, and they matter because they have different
owners. A \textbf{supplier data-entry error} --- where the pack and the submitted
form contradict each other --- goes back to the supplier. A \textbf{resolution
ambiguity} --- where the submission is fine but we cannot tell which of several
similar catalogue products it refers to --- goes to a data steward. A system that
reports both as ``failed'' tells you how many problems you have without telling
you who should fix any of them. The pipeline keeps them separate throughout, and
the audit log records which of the two every held row was.

%============================================================== what was built
\section{What has been built}

Eight stages, each reading and writing a plain file so any stage can be run,
inspected or replaced on its own.

\begin{center}
\begin{tabular}{@{}p{0.6cm}p{4.4cm}p{10.6cm}@{}}
\toprule
\textbf{\#} & \textbf{Stage} & \textbf{What it contributes} \\
\midrule
01 & Vision extraction   & Reads brand, description and barcode off a packaging photograph (Tesseract, two passes with per-field arbitration). \\
02 & Normalisation       & Canonical case, punctuation and units, so \textsf{0.030kg} and \textsf{30g} compare equal. \\
03 & Cross-check         & Compares the three sources pairwise; isolates supplier data-entry errors from resolution ambiguity. \\
04 & Candidate matching  & Ranks the catalogue on four independent signals: fuzzy text, phonetic, TF--IDF character $n$-grams, and barcode. \\
05 & LLM disambiguation  & Called only when retrieval is genuinely undecided; restricted to the shortlist retrieval already produced, and capped so it can never authorise an unreviewed merge. \\
06 & Barcode verification & Three-way agreement across pack, submission and catalogue, with GS1 check-digit validation. \\
07 & Confidence \& routing & Combines the available signals into a calibrated probability and routes: merge / review / reject. \\
08 & Resolve, log, learn & Audit log with full evidence, plus a cache of human-approved corrections so the same fix is never re-derived. \\
\bottomrule
\end{tabular}
\end{center}

Two design decisions are worth surfacing because they are what keep the system
safe rather than merely accurate.

\textbf{Absent evidence is not negative evidence.} A missing barcode used to score
zero --- numerically identical to three barcodes that actively contradict one
another. On a catalogue with no barcodes at all, that capped every row below the
reject threshold and the router emitted a single constant verdict regardless of
input. Confidence is now a weighted mean over the signals actually
\emph{observable} for that row, and missing evidence caps how confident the
system is permitted to be rather than masquerading as evidence against.

\textbf{One signal corroborates nothing.} However strong a text match is, the
pipeline will not merge a record unreviewed unless a second, independent signal
agrees. A perfect text match with nothing to check it against is precisely the
case that merges two different products with similar names. Section~\ref{sec:cal}
quantifies what that rule is worth.

%===================================================================== evidence
\section{Evidence}

Three datasets, deliberately different in what they can and cannot prove.

\begin{center}
\begin{tabular}{@{}p{5.0cm}rrp{5.6cm}@{}}
\toprule
\textbf{Dataset} & \textbf{Catalogue} & \textbf{Top-1} & \textbf{What it establishes} \\
\midrule
Abt--Buy (Leipzig benchmark)     & 1{,}081 & 87.7\% & Text matching against independently written listings. 95\% top-3. \\
Amazon--Google (Leipzig)         & 1{,}363 & 80.2\% & The same engine on a harder, denser catalogue. 95\% top-3. \\
Open Food Facts --- photograph only & 2{,}256 & 24.0\% & Real packaging photo against a real catalogue: what vision alone is worth. \\
\bottomrule
\end{tabular}
\end{center}

The Leipzig figures are the honest text-matching benchmark: both sides are real
retail listings written independently by different merchants, which is the same
problem shape as a supplier submission written independently of our catalogue.

\panel{%
\textbf{A number I am deliberately not quoting.} On the Open Food Facts data the
pipeline scores 100\% when matching the \emph{submission}. That figure is
meaningless: the submission in that dataset is generated by lightly corrupting
the catalogue record, so it is a mangled copy of the answer rather than an
independent description. It is in the evaluation harness to exercise the code
path, not to support a claim. The Leipzig 87.7\% is the number I would defend.}

%========================================================== calibrated confidence
\section{Calibrated confidence}\label{sec:cal}

Routing thresholds are only meaningful if the confidence attached to a decision
is a real probability. Originally it was not: confidence was a raw match margin
divided by a constant chosen against a 30-product test catalogue. On a
1{,}081-product catalogue the median margin for a \emph{correct} match is 0.078,
so almost every row was squashed into the bottom of the range.

The mapping is now fitted on labelled pairs and evaluated on data it never saw:

\begin{center}
\begin{tabular}{@{}lrrrr@{}}
\toprule
& \textbf{Brier} $\downarrow$ & \textbf{Log loss} $\downarrow$ & \textbf{ECE} $\downarrow$ & \textbf{AUC} \\
\midrule
Previous (constant divisor) & 0.530 & 1.544 & 0.656 & 0.908 \\
Calibrated (Platt scaling)  & \textbf{0.076} & \textbf{0.247} & \textbf{0.036} & 0.908 \\
\bottomrule
\end{tabular}
\end{center}

Expected calibration error --- the average gap between the confidence claimed and
the share that turn out correct --- falls from 0.656 to 0.036. Discrimination
(AUC) is unchanged by construction: calibration is monotone, so it cannot alter
which candidate ranks first, only what the number attached to it means.

\textbf{Operationally}, on held-out data this moves \textbf{425 correct matches out of
the reject pile and into the review queue}. Rejection now means ``we have reason
to believe this does not match'' rather than ``we could not tell''.

\panel{%
\textbf{What the corroboration rule is worth, in records.} Disable it and allow the
calibrated 90\% band to merge unreviewed, and 398 submissions auto-merge at 96.7\%
precision --- \textbf{13 wrong records written into the master catalogue}. The band
is behaving exactly as specified; 90\% is simply not a safe bar for an unreviewed
write. This is the kind of trade-off the business should set explicitly, and
calibration is what makes it visible: an uncalibrated score cannot tell you what
your threshold costs.}

%================================================================ vision finding
\section{What packaging photographs are actually worth}

I tested vision extraction on 140 real contributor photographs against the
2{,}256-product catalogue. A single front-of-pack photograph resolves the correct
product \textbf{24\% of the time}.

The limit is not image quality or OCR tuning. Brand recall is only 33\% even on
images where the reader extracted more than 200 characters, because the brand is
frequently a \emph{logo} rather than text --- one Marks \& Spencer pack front reads
``\textsf{FAIRTRADE EST. 1884 SINGLE ORIGIN PERU}'' and never states the brand in
readable type. Separately, the barcode is not on the front of a pack at all, so
it cannot be recovered from a front photograph under any amount of processing.

\textbf{The operational conclusion is to design intake around the scanned barcode
and the typed submission, and treat the photograph as corroboration only.} That
is how goods-in already works --- the barcode comes off a scanner, not a camera ---
and it is what the pipeline's routing rules already assume.

Barcode handling was widened as a result. A real catalogue is 86.8\% EAN-13 but
12.9\% EAN-8 --- small packs where a 13-digit symbol does not physically fit.
Validating EAN-13 only discarded the barcode evidence for one product in seven
and reported it as a \emph{missing} barcode rather than one we had refused to
parse. Accepting all GS1 retail formats took usable barcodes from 85.5\% to
99.2\%.

%============================================================ trust / methodology
\section{Why these numbers can be trusted}

An earlier version of this work reported 100\% accuracy on real data. That was a
\textbf{ground-truth leak}: the evaluation harness copied the catalogue record's own
barcode into the supplier submission, so the matcher was handed the answer on
every row. With the leak closed, the same engine scores 87.7\%.

The leak was worth finding for a reason beyond the corrected figure. It made
every metric look \emph{better}, kept the entire test suite green, and made the
dashboard look healthier. Nothing about it announced itself. The only thing that
surfaced it was two independent measurements of the same quantity disagreeing
--- the pipeline claiming 100\% where a separate code path scored the same engine
at 85\%.

The evaluation harness now carries explicit invariants against that class of
failure, including a check that fails the build if any input field is a
byte-identical copy of a field the input should not have been able to observe.
There are \textbf{63 automated tests}, all runnable offline with no credentials.

%===================================================================== proposal
\section{Proposal}

A six-week pilot on one real supplier feed, scoped to answer a single question:
\emph{what share of submissions can be resolved without human review, at a
wrong-merge rate we are prepared to accept?}

\begin{center}
\begin{tabular}{@{}p{2.0cm}p{7.4cm}p{6.6cm}@{}}
\toprule
\textbf{Weeks} & \textbf{Activity} & \textbf{Output} \\
\midrule
1--2 & Point the pipeline at a real catalogue extract and one live supplier feed. & Baseline accuracy on our own data. \\
3--4 & Data steward labels outcomes on a sample; calibrate confidence against those labels. & Thresholds that mean what they say \emph{here}. \\
5--6 & Shadow-run alongside the current manual process; compare decisions. & Measured automation rate and wrong-merge rate. \\
\bottomrule
\end{tabular}
\end{center}

\textbf{What I need:} a catalogue extract, one supplier feed, and roughly two days
a week of a data steward's time for labelling. \textbf{Nothing is written to any
production system} --- the pilot shadows the existing process and produces a
comparison, not a migration.

Calibration must be fitted on our own catalogue rather than transferred. I
measured this directly: a calibration fitted on one catalogue and applied to
another is systematically overconfident, claiming 65\% where the true rate is
49\%. That is why the steward labelling in weeks 3--4 is not optional --- it is
what makes the confidence figures valid for us.

%======================================================================== risks
\section{Risks and open questions}

\begin{itemize}
  \item \textbf{No-match submissions are untested.} Every labelled pair in the
        public benchmarks is a genuine match, so the reject threshold has never
        been exercised against a submission for a product we simply do not
        stock --- which a real feed contains constantly. This is the first thing
        the pilot should measure.
  \item \textbf{The supplier feed in testing is simulated.} It is derived from
        catalogue records, so it is easier than a real one. The Leipzig figures
        are the conservative estimate; our own feed may sit below them.
  \item \textbf{Phonetic matching assumes English.} It degrades on brands that
        are not, which matters for imported ranges.
  \item \textbf{Automation rate is unknown for our data.} I can evidence that
        the engine works and that its confidence is honest. What share of
        \emph{our} submissions clear an acceptable threshold is exactly what the
        pilot exists to find out, and I would rather establish it than estimate
        it here.
\end{itemize}

\vspace{0.6em}
{\color{rulegrey}\hrule height 0.6pt}
\vspace{0.4em}
{\small\color{black!55}
Working prototype, evaluation harness and 63 automated tests:
\href{https://github.com/jbethell10/MasterData_EntityResolution}%
{github.com/jbethell10/MasterData\_EntityResolution}.
Public data used under ODbL (Open Food Facts) and the Leipzig entity-resolution
benchmark terms.}

\end{document}
